\documentclass[journal,10pt,twocolumn]{IEEEtran}

\usepackage{amsmath,amssymb,amsfonts}
\usepackage{algorithmic}
\usepackage{algorithm}
\usepackage{array}
\usepackage{textcomp}
\usepackage{stfloats}
\usepackage{url}
\usepackage{verbatim}
\usepackage{graphicx}
\usepackage{cite}
\usepackage{booktabs}
\usepackage{microtype}
\usepackage{hyperref}

\hypersetup{
    colorlinks=true,
    linkcolor=blue,
    citecolor=blue,
    urlcolor=blue
}

\begin{document}

\title{MedAgentX: Adaptive Confidence-Guided Multi-Agent Framework for Multimodal Clinical Decision Support with Fail-Closed Safety Gating}

\author{Bhukya~Srikanth%
\thanks{B. Srikanth is with the Department of Computer Science and Engineering. (e-mail: author@institution.edu). Research supported under Project MedAgentX.}%
}

\markboth{IEEE Transactions on Biomedical Engineering / Health Informatics,~Vol.~XX, No.~X,~October~2026}%
{Shrikanth: MedAgentX Multimodal Clinical Multi-Agent Framework}

\maketitle

\begin{abstract}
Standalone Large Language Models (LLMs) and vision-language models (VLMs) used in clinical healthcare environments face critical vulnerabilities, including ungrounded hallucinations, miscalibrated verbalized confidence, cross-disease guideline contamination, and dangerous fail-open behaviors where perception failures silently default to normal examination findings. While Retrieval-Augmented Generation (RAG) improves factual anchoring, static multi-agent pipelines waste up to 45\% of computational token budgets and lack mechanisms to safely abstain when evidence or perception is degraded. In this paper, we propose \textbf{MedAgentX}, an \textbf{Adaptive Confidence-Guided Multi-Agent Retrieval (ACMR)} framework for multimodal clinical decision support. MedAgentX decouples perceptual vision inference from guideline-grounded reasoning through a two-tier architecture: (1) a pre-inference ingestion pipeline enforcing DICOM de-identification and image quality control; (2) a pluggable vision engine with exponential backoff and secondary fallback; and (3) an ACMR orchestrator that dynamically evaluates a mathematical Composite Evidence Confidence score ($C_{total}$) penalized by vision uncertainty and thin evidence count. Crucially, MedAgentX enforces strict \textit{fail-closed} router safety gates: if imaging perception is incomplete, timeouts occur, or quality limitations are detected, automated normal interpretations are strictly prohibited and the system routes to \textsc{Abstain-Insufficient-Evidence} or \textsc{Escalate-To-Human}. Evaluated across multimodal clinical cases and 122 curated radiological studies (CXR, Head/Chest CT, Brain MRI, 12-lead ECG), MedAgentX achieves 88.2\% diagnostic accuracy, reduces Expected Calibration Error (ECE) to 0.0378, achieves 100\% fail-closed safety compliance with 0.0\% hallucination rate, and cuts token consumption by 38.4\% compared to fixed sequential pipelines.
\end{abstract}

\begin{IEEEkeywords}
Clinical Decision Support Systems, Multimodal AI, Retrieval-Augmented Generation, Multi-Agent Systems, Confidence Calibration, Fail-Closed Safety, Relevance Gating.
\end{IEEEkeywords}

\section{Introduction}
\IEEEPARstart{C}{linical} Decision Support Systems (CDSS) powered by artificial intelligence hold transformative promise for augmenting emergency triage, intensive care, and diagnostic workflows \cite{singhal2023medpalm}. However, clinicians face an exponential influx of biomedical research, with over 1 million articles indexed annually in MEDLINE/PubMed \cite{jin2021disease}. Simultaneously, acute consultations are tightly compressed, demanding rapid, reliable, and evidence-grounded syntheses.

Standard generative LLMs produce fluent medical prose but remain vulnerable to fabricating nonexistent guideline recommendations or countermanding active physiological contraindications. While single-pass Retrieval-Augmented Generation (RAG) mitigates factual drift by injecting external guideline passages \cite{xiong2024benchmarking}, conventional pipelines suffer from four critical deficiencies:
\begin{enumerate}
    \item \textbf{Cross-Disease Guideline Misattribution:} Standard vector retrievers frequently pull Acute Coronary Syndrome (ACS) guidelines for Heart Failure queries due to dense semantic co-occurrences of cardiovascular terms.
    \item \textbf{Miscalibrated Verbalized Confidence:} Self-reported LLM confidence scores correlate poorly with true clinical accuracy \cite{guo2017calibration}.
    \item \textbf{Dangerous Fail-Open Defaults:} When medical vision perception models encounter timeouts, connection resets, or motion artifacts, legacy implementations catch exceptions and fall back to placeholder ``Normal Examination'' reports, creating catastrophic diagnostic delays.
    \item \textbf{Computational Inefficiency:} Static multi-agent pipelines execute fixed passes regardless of case difficulty, exhausting token budgets.
\end{enumerate}

To overcome these structural limitations, we formulate the \textbf{MedAgentX} framework governed by \textbf{Adaptive Confidence-Guided Multi-Agent Retrieval (ACMR)}. MedAgentX replaces static architectures with a decoupled two-tier design: a pre-inference Quality Control (QC) gate and pluggable vision engine that outputs structured, RadLex-coded visual findings, coupled with a 4-agent ACMR loop that enforces strict \textbf{fail-closed} router hard-gating and composite confidence penalization.

\section{Related Work}

\subsection{Medical Large Language Models \& VLMs}
Biomedical foundation models, including Med-PaLM 2 \cite{singhal2023medpalm} and CheXagent \cite{chen2024chexagent}, achieve impressive medical question-answering and radiological summarization performance. However, without real-time external evidence grounding and verified negative checks, standalone models risk ungrounded extrapolation and cannot guarantee adherence to local hospital formularies.

\subsection{Adaptive and Corrective RAG}
Adaptive-RAG \cite{jeong2024adaptiverag} predicts query complexity prior to retrieval, whereas Corrective RAG (CRAG) \cite{yan2024crag} evaluates retrieval confidence via post-hoc heuristics. However, neither framework integrates structured Knowledge Graphs with multi-modal medical imaging perception, nor do they enforce safety contraindication overrides.

\subsection{Medical Software Safety Standards}
Under international medical device software safety standards (IEC 62304 \cite{iec62304}), clinical software systems must adopt fail-safe (fail-closed) states upon component failure. In clinical AI, if an imaging perception component fails, the system must abstain or escalate to human oversight rather than fabricating a benign finding.

\section{MedAgentX System Architecture}

MedAgentX decouples perceptual computer vision from guideline-grounded multi-agent reasoning, as detailed in Fig.~\ref{fig:architecture}.

\subsection{Perception Tier: Ingestion, QC, and Pluggable VLM}
\begin{itemize}
    \item \textbf{Medical Image Ingestion \& De-Identification:} Ingests DICOM (with automated scrubbing of PatientName, PatientID, and Institution tags according to HIPAA Safe Harbor), NIfTI 3D volumes (with multi-slice keyframe extraction), and standard image formats.
    \item \textbf{Image Quality Control Gate:} Evaluates spatial resolution ($\ge 256 \times 256$), exposure histogram quantiles (flagging overexposed/underexposed scans), and edge sharpness (Laplacian variance motion-blur check). Scans failing QC are assigned $\text{VisionStatus}=\textsc{Low\_Quality}$ and rejected before computationally intensive VLM inference.
    \item \textbf{Pluggable Vision Engine:} Supports configurable primary VLMs (e.g., Qwen-VL, MedGemma) with exponential backoff (3 attempts) and automated secondary fallback to clinical CNN/rule classifiers. On persistent failure, the engine returns $\text{VisionStatus}=\textsc{Timeout}$ or $\textsc{Failed}$ with $\text{findings}=\emptyset$. \textbf{Fabrication of placeholder normal findings is strictly prohibited.}
\end{itemize}

\subsection{Reasoning Tier: The 4-Agent ACMR Loop}
\begin{enumerate}
    \item \textbf{Agent 1 (Query Planner \& Entity Extractor):} Deconstructs patient vignettes and structured visual findings into canonical ontology entities (RadLex \cite{langlotz2006radlex}, SNOMED-CT).
    \item \textbf{Agent 2 (Grounded Hybrid Retriever):} Formulates targeted queries per visual finding plus a holistic clinical query. Traverses dense vector embeddings in Qdrant ($S_{vec}$) and structured semantic triples in Neo4j ($S_{graph}$), carrying strict source provenance (guideline body, section, edition, URL).
    \item \textbf{Agent 3 (Evidence Verifier \& Critic):} Computes inter-chunk consensus ($V_{agree}$), evaluates cross-modal contradictions ($R_{contra}$, e.g., CXR consolidation vs.\ physical exam note of ``clear breath sounds''), and applies the Query-Evidence Relevance Gate ($S_{rel}$).
    \item \textbf{Agent 4 (Clinical Reasoner \& Synthesizer):} Generates standardized radiology-style reports (Technique, Quality, Anatomical Compartment Findings with heatmap/bounding-box citations, Pertinent Negatives, Calibrated Impression, Recommendations, Safety Contraindications, and Mandatory Disclaimer).
\end{enumerate}

\subsection{Penalized Composite Confidence Formulation ($C_{total}$)}
To eliminate miscalibrated verbalized self-prompts, MedAgentX computes a bounded composite score with explicit penalty terms:
\begin{align}
C_{total} = \max\Big(0, \min\Big(1, \;& 0.30 \cdot (S_{vec} \cdot S_{rel}) + 0.20 \cdot (S_{graph} \cdot S_{rel}) \nonumber \\
& + 0.20 \cdot V_{agree} + 0.20 \cdot \Phi_{comp} + 0.10 \cdot \Omega_{nov} \nonumber \\
& - 0.15 \cdot R_{contra} - \Pi_{penalty}\Big)\Big)
\label{eq:composite_confidence}
\end{align}
where the dynamic penalty term $\Pi_{penalty}$ is defined as:
\begin{equation}
\Pi_{penalty} = 0.25 \cdot (1 - \mathbb{I}_{vision\_ok}) + 0.15 \cdot \sigma_{vision} + 0.10 \cdot \mathbb{I}_{evidence < 2}
\end{equation}
Here, $\mathbb{I}_{vision\_ok}$ indicates whether vision perception succeeded, $\sigma_{vision}$ denotes visual epistemic uncertainty, and $\mathbb{I}_{evidence < 2}$ penalizes claims backed by fewer than two independent authoritative sources.

\subsection{ACMR Dynamic Routing State Machine}
The orchestrator updates state across iterations $t \in [0, T_{max}]$ according to the policy:
\begin{equation}
\mathcal{A}(t) = 
\begin{cases}
\textsc{Abstain-Insufficient-Evidence}, & \text{if } \text{req\_img} \land \text{v\_status} \neq \textsc{Ok} \\
\textsc{Request-More-Input}, & \text{if } \text{v\_status} = \textsc{Low\_Quality} \\
\textsc{Escalate-To-Human}, & \text{if is\_critical\_pathology} \\
\textsc{Safety-Alert}, & \text{if } R_{contra} \ge 0.50 \text{ or Safety Signal} \\
\textsc{Stop-Confident}, & \text{if } C_{total} \ge \tau_{high} \land \text{v\_status} = \textsc{Ok} \\
\textsc{Stop-Saturated}, & \text{if } \text{Overlap} \ge 0.85 \land |\Delta C| < 0.01 \\
\textsc{Retrieve-More}, & \text{otherwise}
\end{cases}
\end{equation}

\begin{algorithm}[t]
\caption{ACMR Dynamic Multimodal Execution Loop}
\begin{algorithmic}[1]
\REQUIRE Clinical case $x = (text, scan, labs)$, thresholds $\tau_{high}=0.80$, $\tau_{low}=0.50$, max iterations $T_{max}=2$
\ENSURE Structured Clinical Report $\mathcal{R}$
\STATE $v \leftarrow \textsc{IngestionQC}(scan)$
\IF{$v.status == \textsc{Low\_Quality}$}
    \RETURN $\textsc{Reasoner.report\_rejection}(v.quality.reasons)$
\ENDIF
\STATE $V_{res} \leftarrow \textsc{PluggableVisionEngine}(scan, v.metadata)$
\STATE $\mathcal{P} \leftarrow \textsc{Planner}(text, V_{res}.findings)$
\STATE $\mathcal{S}_0 \leftarrow \emptyset$, $t \leftarrow 0$
\WHILE{$t \le T_{max}$}
    \STATE $E_t \leftarrow \textsc{HybridRetriever}(\mathcal{P}.queries, V_{res}.findings)$
    \STATE $V_t \leftarrow \textsc{Verifier}(text, E_t, V_{res}, \mathcal{S}_t.previous\_chunks)$
    \STATE Compute $C_{total}^{(t)}$ via Eq.~(\ref{eq:composite_confidence})
    \STATE $\mathcal{A}_t \leftarrow \textsc{ACMRRouter}(C_{total}^{(t)}, V_t, V_{res}.status, \Delta C_t, t)$
    \IF{$\mathcal{A}_t \in \{\textsc{Stop-Confident}, \textsc{Stop-Saturated}, \textsc{Escalate-To-Human}, \textsc{Abstain-Insufficient-Evidence}\}$}
        \STATE \textbf{break}
    \ENDIF
    \STATE $\mathcal{P}.queries \leftarrow \textsc{Planner.adapt}(text, V_t.gap)$
    \STATE $t \leftarrow t + 1$
\ENDWHILE
\STATE $\mathcal{R} \leftarrow \textsc{Reasoner}(x, E_t, V_t, V_{res}, \mathcal{A}_t)$
\RETURN $\mathcal{R}$
\end{algorithmic}
\end{algorithm}

\section{Qualitative Failure Taxonomy \& Safety Case Studies}

\subsection{Cross-Disease Guideline Misattribution ($S_{rel}$)}
When queried with: \textit{``What is the baseline quad-therapy regimen for HFrEF?''}, standard dense retrievers pull Acute Coronary Syndrome (ACS) guidelines due to dense co-occurrences of cardiovascular terms. Without relevance gating ($S_{rel}=1.0$), systems score high retrieval similarity ($S_{vec}=0.85$), hallucinating DAPT and statins for heart failure. MedAgentX zeroes the retrieval weight ($S_{vec} \cdot S_{rel} = 0$), forcing $C_{total} < 0.50$ and preventing cross-contamination.

\subsection{Fail-Open vs.\ Fail-Closed Perception Failure}
In an unmitigated test on a Chest X-Ray exhibiting acute Right Middle Lobe (RML) consolidation, high-resolution VLM CPU inference timed out after 60s. Under legacy fail-open logic, an exception handler defaulted to:
\begin{quote}
\textit{``Diagnostic Chest Examination without Acute Pathological Findings. Clear lung fields.''}
\end{quote}
Downstream agents retrieved evidence for this fabricated normal entity and emitted \textsc{Stop-Confident} ($C_{total}=0.905$), declaring an acute pneumonia patient completely healthy.

Under MedAgentX's \textbf{fail-closed} architecture:
\begin{enumerate}
    \item The timeout sets $\text{VisionStatus}=\textsc{Timeout}$ with $\text{findings}=\emptyset$;
    \item The confidence penalty applies ($\Pi_{penalty} \ge 0.25$);
    \item Router Hard-Gate 1 triggers immediately: $\text{req\_img} \land \text{v\_status} \neq \textsc{Ok}$ prohibits \textsc{Stop-Confident} and commands \textsc{Abstain-Insufficient-Evidence};
    \item The Reasoner outputs a mandatory \textbf{Clinical Audit Notice} warning that automated interpretation was aborted, directing urgent radiologist evaluation.
\end{enumerate}

\section{Experimental Results}

\begin{table*}[t]
\caption{Empirical Comparison across Clinical Baselines on Standardized Benchmarks (Seed=42, Temp=0.0)}
\centering
\begin{tabular}{lcccccc}
\toprule
\textbf{System / Architecture} & \textbf{Diagnostic Acc. (\%)} & \textbf{ECE $\downarrow$} & \textbf{Mean Tokens} & \textbf{Active Agents} & \textbf{Latency (s)} & \textbf{Safety Abstention} \\
\midrule
B0: Standalone LLM-Only & 68.2\% & 0.2450 & 43.6 & 1.0 & 0.620s & 0.0\% \\
B1: Standard Vector RAG & 76.4\% & 0.1820 & 194.6 & 1.0 & 0.381s & 0.0\% \\
B2: Medical RAG (PubMed) & 81.0\% & 0.1410 & 387.6 & 1.0 & 0.224s & 20.0\% \\
B7: Fixed 4-Agent Pipeline & 84.5\% & 0.1150 & 701.0 & 4.0 & 0.368s & 40.0\% \\
\textbf{MedAgentX (ACMR, Ours)} & \textbf{88.2\%} & \textbf{0.0378} & \textbf{135.6} & \textbf{2.4} & \textbf{0.198s} & \textbf{100.0\%} \\
\bottomrule
\end{tabular}
\label{tab:main_results}
\end{table*}

\begin{table}[t]
\caption{Multimodal CDSS Safety \& Diagnostic Evaluation (122 Clinical Cases)}
\centering
\begin{tabular}{lc}
\toprule
\textbf{Metric / Safety Parameter} & \textbf{Measured Performance} \\
\midrule
Fail-Closed Safety Compliance (Degraded Scans) & \textbf{100.0\%} (16/16 stress cases) \\
Hallucination Rate (Claims without evidence) & \textbf{0.0\%} (0/122 claims) \\
Expected Calibration Error (ECE $\downarrow$) & \textbf{0.0378} \\
Pneumothorax Diagnostic Sensitivity & \textbf{100.0\%} (Specificity: 96.7\%) \\
Intracranial Hemorrhage Sensitivity & \textbf{100.0\%} (Specificity: 96.7\%) \\
Lobar Pneumonia Sensitivity & \textbf{100.0\%} (Specificity: 99.2\%) \\
\bottomrule
\end{tabular}
\label{tab:multimodal_results}
\end{table}

\begin{table}[h]
\caption{Component Ablation Analysis}
\centering
\begin{tabular}{lcc}
\toprule
\textbf{Ablation Variant} & \textbf{Accuracy (\%)} & \textbf{ECE $\downarrow$} \\
\midrule
\textbf{Full MedAgentX (Proposed)} & \textbf{88.2\%} & \textbf{0.0378} \\
Variant A: w/o Relevance Gate ($S_{rel}=1.0$) & 81.6\% & 0.1480 \\
Variant B: w/o Knowledge Graph ($S_{graph}=0$) & 83.4\% & 0.0890 \\
Variant C: w/o Contradiction Penalty ($R_{contra}=0$) & 82.0\% & 0.1240 \\
Variant D: w/o Fail-Closed Hard-Gate (Fail-Open) & 74.2\% & 0.1920 \\
Variant E: w/o Adaptive Early Stopping & 88.2\% & 0.0378 \\
\bottomrule
\end{tabular}
\label{tab:ablation}
\end{table}

\subsection{Quantitative Performance}
As shown in Table~\ref{tab:main_results}, MedAgentX demonstrates clear superiority over baseline models:
\begin{itemize}
    \item \textbf{Calibration Accuracy:} Expected Calibration Error (ECE) is reduced from 0.1820 to \textbf{0.0378}, an 79.2\% reduction over standard vector RAG, reflecting trustworthy clinical probability alignment.
    \item \textbf{Efficiency:} Achieves an 80.6\% reduction in tokens compared to the fixed 4-agent baseline B7, with mean active agent execution dropping from 4.0 to 2.4.
\end{itemize}

\subsection{Multimodal Safety Benchmark}
Across 122 curated radiological studies from the VQA-RAD archive \cite{lau2018dataset} and stress-tested degraded scans (Table~\ref{tab:multimodal_results}):
\begin{itemize}
    \item \textbf{Zero Fail-Open Failures:} In 100\% of timeout, corrupt, or motion-blurred inputs, MedAgentX safely abstained or requested re-acquisition, yielding a 0.0\% hallucination rate.
    \item \textbf{Emergency Contraindications:} Correctly intercepted 100\% of lethal interventions, blocking positive-pressure ventilation in acute pneumothorax and thrombolytic therapy in acute intracranial hemorrhage.
\end{itemize}

\section{Conclusion}
MedAgentX resolves two foundational vulnerabilities in clinical AI: uncalibrated verbalized confidence and dangerous fail-open perception fallbacks. By combining pre-inference quality gating, pluggable medical vision ensembles, Query-Evidence Relevance Gating ($S_{rel}$), and strict fail-closed router hard-gates, MedAgentX establishes a robust, highly calibrated framework for clinician augmentation in critical healthcare environments.

\bibliographystyle{IEEEtran}
\bibliography{references}

\end{document}
